% Research Methods Section
% PhD Student 3 — Machine learning models, pipeline, and applied tests
% Paper: China's out-of-home food consumption
% No \begin{document}; for inclusion in main manuscript.

\section{Research Methods}
\label{sec:methods}

\subsection{Prediction Target and Outcome Variable}

The training and prediction target is the \emph{at-home consumption ratio} (ratio of at-home to total at-home plus out-of-home consumption), defined as
\begin{equation}
  \text{ratio} = \frac{\text{at-home}}{\text{at-home} + \text{out-of-home}} \in (0,\,1].
\end{equation}
The \emph{out-of-home consumption coefficient} is then defined as the reciprocal of this ratio:
\begin{equation}
  \text{out-of-home coefficient} = \frac{1}{\text{ratio}} = \frac{\text{at-home} + \text{out-of-home}}{\text{at-home}}.
\end{equation}
Total (at-home plus out-of-home) consumption is obtained in post-processing as
\begin{equation}
  \text{at-home} + \text{out-of-home} = \text{at-home} \times \text{out-of-home coefficient}.
\end{equation}
This relation is applied in the post-processing step (\texttt{postprocess\_predictions.py}) to derive provincial and national consumption quantities from predicted out-of-home coefficients and at-home consumption from \texttt{data\_q.csv}.

For interpretation, the coefficient is a \emph{scaling factor} that maps at-home quantities into total quantities under aligned definitions. The implied out-of-home share of total consumption is
\begin{equation}
  \text{share}_{\text{out}} = \frac{\text{out-of-home}}{\text{total}} = 1 - \frac{1}{\text{out-of-home coefficient}}.
\end{equation}
We emphasize that this paper estimates scaling factors for accounting and projection, not behavioral parameters or causal effects.

Predictions are produced for \textbf{15 food categories}: four grain types---paddy rice (\texttt{daogu}), wheat (\texttt{xiaomai}), beans (\texttt{doulei}), and coarse grains (\texttt{zaliang})---plus vegetables (\texttt{shucai}); meat (pork \texttt{zhurou}, beef \texttt{niurou}, mutton \texttt{yangrou}, poultry \texttt{qinlei}); aquatic products (\texttt{shuichanpin}); eggs (\texttt{danlei}); dairy (\texttt{nailei}); fruit (\texttt{guaguo}); edible oil (\texttt{youliao}); and sugar (\texttt{tang}).

\subsection{Features and Data Preparation}

Features are constructed in \texttt{data\_preparation\_advanced.py} in the function \texttt{load\_and\_prepare\_data\_advanced}. The feature vector used for both training and prediction is
\begin{equation}
  X = [\texttt{indinc},\, \texttt{income},\, \texttt{urbanrate},\, \texttt{oldrate},\, \texttt{engel},\, \texttt{foodpindex},\, \texttt{zyywsr},\, \texttt{cyyye},\, \texttt{cyyysr},\, \text{production variables},\, \texttt{T1},\, \texttt{wave}].
\end{equation}

\paragraph{Micro-level.}
\texttt{indinc} is household income aligned with the macro prediction table. When income-distribution matching is enabled, it is generated by \texttt{match\_income\_copula} so that the (rank-preserving) distribution of income in the macro scenario matches the micro survey distribution after scaling to macro means (see Section~\ref{sec:copula}).

\paragraph{Macro-level (province $\times$ year).}
Macro features include \texttt{income}, \texttt{urbanrate}, \texttt{oldrate}, \texttt{engel}, \texttt{foodpindex}, \texttt{zyywsr}, \texttt{cyyye}, \texttt{cyyysr}, and production variables. Production variables are merged from \texttt{data\_production1.csv}, \texttt{data\_production2.csv}, and \texttt{data\_production3.csv} by \texttt{wave} and \texttt{T1}; they include \texttt{qliangshi}, \texttt{qdaogu}, \texttt{qxiaomai}, \texttt{qyumi}, \texttt{qdoulei}, \texttt{qshulei}, \texttt{qyouliao}, \texttt{qshuiguo}, \texttt{qroulei}, \texttt{qzhurou}, \texttt{qniurou}, \texttt{qyangrou}, \texttt{qnailei}, \texttt{qqindan}, and \texttt{qshuichan}.

\paragraph{Identifiers.}
\texttt{T1} is the province code and \texttt{wave} is the year. Prediction is performed only for \texttt{wave} $\geq$ 2015 (\texttt{MIN\_PREDICT\_YEAR}).

Training uses the micro dataset \texttt{data.csv} with the above features and category-specific \texttt{ratio\_*} (at-home over at-home plus out-of-home) as the target, without truncation. Prediction uses the macro table \texttt{data2012.csv} and production tables to build the province$\times$year feature matrix for each category.

\subsection{Machine Learning Models}
\label{sec:models}

We use \textbf{12} machine learning models for coefficient prediction and, where applicable, for imputation. Each is implemented in a dedicated script (\texttt{predict\_<model>.py} or \texttt{predict\_fttransformer\_advanced.py}).

\paragraph{Tree-based and linear.}
\begin{itemize}
  \item \textbf{LightGBM}: gradient boosting (LGBMRegressor); fixed hyperparameters; used for coefficient prediction and in imputation selection.
  \item \textbf{XGBoost}: extreme gradient boosting; time-series cross-validation (TimeSeriesSplit, 3 splits) in training; coefficient prediction and imputation.
  \item \textbf{CatBoost}: gradient boosting with categorical support; TimeSeriesSplit (3 splits); coefficient prediction and imputation.
  \item \textbf{Random Forest}: ensemble of trees; coefficient prediction and imputation.
  \item \textbf{Lasso}: L1-regularized linear regression; coefficient prediction and imputation.
  \item \textbf{Linear (Ridge)}: L2-regularized linear regression; coefficient prediction and imputation.
\end{itemize}

\paragraph{Deep tabular.}
\begin{itemize}
  \item \textbf{MLP}: multilayer perceptron (\texttt{sklearn.neural\_network.MLPRegressor}); early stopping; fallback for TabNet/TabPFN when unavailable; coefficient prediction and imputation.
  \item \textbf{TabNet}: attention-based tabular model (pytorch-tabnet TabNetRegressor); \texttt{patience} for early stopping; fallback to MLP if unavailable; coefficient prediction and imputation.
  \item \textbf{TabM}: parameter-efficient MLP-style model (yandex-research/tabm); current implementation may fall back to MLP; coefficient prediction and imputation.
  \item \textbf{TabPFN}: prior-fitted network (TabPFNRegressor); local checkpoint when available; fallback to MLP; coefficient prediction and imputation.
  \item \textbf{ResNet (tabular)}: residual network for tabular data; coefficient prediction and imputation.
  \item \textbf{FT-Transformer}: feature-tokenizer Transformer (PyTorch); \texttt{DirectRatioOptimizer} with Optuna for coefficient prediction; dedicated \texttt{FTRegressionModel} for imputation (home/total, no sigmoid). FT-Transformer uses its own imputation when selected as best; it is still evaluated in imputation selection via \texttt{evaluate\_imputation\_cv\_fttransformer} and can be chosen as the single best model.
\end{itemize}

All 12 models participate in the imputation model selection (\texttt{run\_imputation\_selection.py}); the best model (e.g., FT-Transformer, as in \texttt{data/imputation\_selection\_results.csv}) is used to produce the unified imputed sample \texttt{data/imputed\_ratios\_best.csv}.

\subsection{Pipeline Overview}
\label{sec:pipeline}

The workflow consists of the following steps.

\begin{enumerate}
  \item \textbf{Data preparation.}
  Load micro \texttt{data.csv}, macro \texttt{data2012.csv}, and production tables; merge production by \texttt{wave} and \texttt{T1}; restrict prediction to years $\geq$ 2015. Optionally apply income-distribution matching (rank-preserving scaling) to obtain \texttt{indinc} for the macro table. Output: prepared micro data, prediction feature matrix \texttt{X\_pred\_feats}, feature list \texttt{feature\_cols}, category map, and set of known provinces.

  \item \textbf{Imputation model selection} (\texttt{run\_imputation\_selection.py}).
  For each category, the 12 models (including FT-Transformer) are evaluated on \emph{observed} at-home/total data using 5-fold cross-validation. Each model predicts at-home and total consumption; the ratio is computed and compared to the observed ratio. Metrics (MAE, RMSE, R²) are averaged across categories per model; the model with the lowest average MAE is selected. The selected model is then used to impute missing at-home/total (and thus ratio) for all categories; results are saved as \texttt{data/imputed\_ratios\_best.csv} and the model name in \texttt{data/best\_imputation\_model.txt}. Aggregated metrics and ranks are written to \texttt{data/imputation\_selection\_results.csv}.

  \item \textbf{Coefficient prediction (main pipeline).}
  If \texttt{data/imputed\_ratios\_best.csv} exists, each model uses the unified imputed sample (\texttt{ratio\_filled\_*}) to train the coefficient model and to predict the at-home consumption ratio for the macro grid. Otherwise, each model imputes internally. For each category, the out-of-home coefficient is obtained as the reciprocal of the predicted ratio. Missing provinces are filled via spatial interpolation (inverse-distance weighting, IDW). Outputs: \texttt{predictions\_<model>.csv} (out-of-home coefficients by province and year).

  \item \textbf{Grain structure.}
  Four grain shares (paddy, wheat, beans, coarse grains) are predicted from micro data (e.g., paddy from total rice/0.7, then shares of total grain mass). For provinces missing in the micro data, shares are obtained by spatial interpolation (IDW) combined with provincial production shares. Results are saved as \texttt{grain\_structure\_predictions\_<model>.csv} and used in post-processing.

  \item \textbf{Post-processing} (\texttt{postprocess\_predictions.py}).
  Using predicted out-of-home coefficients, grain structure, \texttt{data\_q.csv} (provincial per-capita at-home consumption), and \texttt{procince\_pop.csv} (population): provincial per-capita at-home consumption for each grain type is \texttt{q\_liangshi} times the corresponding share; provincial per-capita at-home plus out-of-home is per-capita at-home times the out-of-home coefficient. For non-grain categories, provincial per-capita at-home is taken from \texttt{data\_q.csv} and multiplied by the out-of-home coefficient to get per-capita at-home plus out-of-home. National out-of-home coefficients and per-capita consumption are computed as population-weighted averages across provinces. Outputs: \texttt{results\_province\_<model>.csv}, \texttt{results\_national\_<model>.csv}.

  \item \textbf{Robustness (no imputation).}
  Each model is also run without imputation: only samples with observed ratio (and consumption $>0$) are used for training and prediction. Results are saved as \texttt{predictions\_<model>\_robust.csv} and can be compared to the main pipeline.

  \item \textbf{Bootstrap prediction intervals} (\texttt{run\_bootstrap.py}).
  For a given model, prediction is run multiple times (default 30) with different random seeds (\texttt{PREDICT\_SEED}). For each (Province, Year) and each category, the mean and 2.5\% and 97.5\% percentiles are computed and saved as \texttt{predictions\_<model>\_bootstrap.csv} (columns \texttt{*\_Mean}, \texttt{*\_Lower}, \texttt{*\_Upper}).
\end{enumerate}

\subsection{Supporting Methods}

\subsubsection{Income-distribution matching (rank-preserving scaling)}
\label{sec:copula}

Micro data contain household income (\texttt{indinc}), while the macro prediction table has only province--year mean income (\texttt{income}). To align the training and prediction feature spaces, we apply income-distribution matching implemented in \texttt{match\_income\_copula} in \texttt{data\_preparation\_advanced.py}. Despite the function name, the procedure is best interpreted as a \emph{rank-preserving quantile-scaling} approach rather than a fully parametric Copula model: for each province--year in the macro table, the micro income distribution is taken from the empirical distribution of household income in the corresponding micro sample where available (or from the pooled micro sample otherwise). A scale factor is computed so that the mean of the scaled distribution equals the macro mean income for that province--year (scale factor = macro mean / micro mean); empirical quantiles of the micro distribution are then multiplied by this factor to produce synthetic \texttt{indinc} values for the macro table. This step runs once during data preparation and is independent of food category.

\subsubsection{Spatial interpolation for missing provinces}

A province is classified as \emph{missing} if it has no micro data in the survey (i.e., it does not appear in the micro sample at all), as opposed to having no data only for a particular year. For each year, coefficients for missing provinces are filled using spatial interpolation from provinces that do have micro data. \textbf{Out-of-home coefficients:} \texttt{apply\_kriging\_interpolation} in \texttt{data\_preparation\_advanced.py} uses inverse-distance weighting (IDW; distance-based weights) to fill missing provinces for the out-of-home coefficient, by year. \textbf{Grain structure:} \texttt{apply\_grain\_structure\_interpolation} uses IDW (inverse-distance squared weights) combined with the province’s production share of total grain for that type; the final share is a blend of the spatial interpolant and the production-based share. An alternative IDW-only path is available in \texttt{apply\_spatial\_interpolation} for scripts that call it.

\subsubsection{Grain structure (four shares)}

Grain is split into four types: paddy (稻谷), wheat (小麦), beans (豆类), and coarse grains (杂粮). From micro data, paddy mass is derived as total rice/0.7; total grain mass is the sum of the four components. Each share is modeled separately; for provinces without micro data, the share is obtained by IDW plus production share as above. The four shares are normalized to sum to one and written to \texttt{grain\_structure\_predictions\_<model>.csv}. In post-processing, provincial per-capita at-home consumption for each grain type is \texttt{q\_liangshi} times the corresponding share; the out-of-home coefficient for that type is then applied to obtain per-capita at-home plus out-of-home consumption.

\subsection{Evaluation and Regularization}

\begin{itemize}
  \item \textbf{Imputation selection.} 5-fold cross-validation on observed (at-home, total) pairs; metrics MAE, RMSE, R²; selection by lowest average MAE across categories; 12 models compared; results in \texttt{data/imputation\_selection\_results.csv} and \texttt{data/imputation\_evaluation\_detail.csv}.
  \item \textbf{Robustness.} Training and prediction using only observed ratio (no imputation), yielding \texttt{predictions\_<model>\_robust.csv} for comparison with the main pipeline.
  \item \textbf{Bootstrap.} Multiple runs with different seeds; mean and 2.5\%/97.5\% percentiles per (Province, Year) and category in \texttt{predictions\_<model>\_bootstrap.csv}.
  \item \textbf{Cross-validation.} 5-fold CV in imputation selection; TimeSeriesSplit (3 splits) in coefficient training for several models (e.g., XGBoost, CatBoost, FT-Transformer).
  \item \textbf{Early stopping.} Used where applicable: e.g., MLP (\texttt{early\_stopping=True}), TabNet (\texttt{patience}), FT-Transformer (validation loss and \texttt{patience}); reduces overfitting.
  \item \textbf{Regularization.} L1/L2 in Lasso/Ridge; fixed or Optuna-tuned hyperparameters for tree and deep models (FT-Transformer uses Optuna for coefficient prediction).
\end{itemize}

All methodological claims above are aligned with the implementation in \texttt{data\_preparation\_advanced.py}, \texttt{run\_imputation\_selection.py}, \texttt{run\_bootstrap.py}, \texttt{postprocess\_predictions.py}, and the \texttt{predict\_*.py} scripts, and with the project documentation.
